\documentclass[11pt,a4paper,twoside]{article}

\usepackage[a4paper,margin=1in]{geometry}
\usepackage[T1]{fontenc}
\usepackage[utf8]{inputenc}
\usepackage{lmodern}
\usepackage{amsmath,amsthm,amssymb,mathtools}
\usepackage{thmtools}         % before hyperref/cleveref
\usepackage[hidelinks]{hyperref}
\usepackage[nameinlink,capitalize]{cleveref}
\usepackage{bm}

% ===== Theorem styles =====
\theoremstyle{plain}
\newtheorem{theorem}{Theorem}[section]
\newtheorem{prop}[theorem]{Proposition}
\newtheorem{lemma}[theorem]{Lemma}
\newtheorem{cor}[theorem]{Corollary}
\theoremstyle{definition}
\newtheorem{defn}[theorem]{Definition}
\theoremstyle{remark}
\newtheorem{remark}[theorem]{Remark}

% ===== Macros =====
\newcommand{\RR}{\mathbb{R}}
\newcommand{\dd}{\,\mathrm{d}}
\newcommand{\Tr}{\mathrm{Tr}}
\newcommand{\Vol}{\mathrm{Vol}}
\newcommand{\MTT}{\mathrm{MTT}}
\DeclareMathOperator{\Ric}{Ric}
\newcommand{\abs}[1]{\left\lvert #1 \right\rvert}
\newcommand{\brk}[1]{\left( #1 \right)}
\newcommand{\sbrk}[1]{\left[ #1 \right]}
\newcommand{\cbrk}[1]{\left\{ #1 \right\}}
\newcommand{\half}{\tfrac{1}{2}}
\newcommand{\veps}{\varepsilon}
\newcommand{\E}{\mathcal{E}}
\newcommand{\Hcal}{\mathcal{H}}
\newcommand{\Acal}{\mathcal{A}}
\newcommand{\Bcal}{\mathcal{B}}
\newcommand{\Ccal}{\mathcal{C}}
\newcommand{\Dcal}{\mathcal{D}}
\newcommand{\Ocal}{\mathcal{O}}
\newcommand{\Ical}{\mathcal{I}}
\newcommand{\gBI}{\gamma} % Barbero-Immirzi parameter
\newcommand{\epsijk}{\epsilon^{\,ijk}}
\newcommand{\epsabc}{\epsilon^{\,abc}}
%\newcommand{\gBI}{\gamma_{\mathrm{BI}}}

% ===== Title =====
\title{\Large\bf  Modal Triplet Theory: From MTT to Loop Quantum Gravity:\\[0.25em]
A Canonical and Spin--Foam Embedding with Predictive Immirzi Map}
\author{Peter Nero}
\date{September 19, 2025}

\begin{document}
\maketitle

\begin{abstract}
We derive Loop Quantum Gravity (LQG) from the coherent fixed-point sector of Modal Triplet Theory (MTT). Starting with a $3+1$ foliation of the 4D face of the MTT fixed point, we show: (i) the Ashtekar--Barbero connection $A^i_a=\Gamma^i_a+\gBI K^i_a$ and densitized triad $E^a_i$ arise canonically from the modal tetrad/spin connection; (ii) the holonomy--flux $*$-algebra, cylindrical consistency, and the Ashtekar--Lewandowski representation follow from the $\Pi$-projected symmetry and yield the LOST uniqueness properties; (iii) the Gauss, diffeomorphism, and Hamiltonian constraints descend from MTT gauge/diffeomorphism invariance and from the fixed-point Hamiltonian, including a Master-constraint form; (iv) area/volume spectra of LQG are recovered with a \emph{computed} Barbero--Immirzi parameter $\gBI=\gBI_{\mathrm{MTT}}(\Theta)$ determined by modal overlaps/bottlenecks; and (v) spin--foam dynamics (EPRL/FK) arise from a Holst-like MTT effective action via BF + simplicity constraints, with large-spin Regge asymptotics. The embedding is consistent with the UV-finite, causal perturbative QG sector previously derived from MTT (Stieltjes/Bernstein two-point function and SPT Gaussian damping), and identifies cross-checks linking $\gBI$, black-hole entropy, and cosmological loop effects.
\end{abstract}

\tableofcontents
\newpage

% =========================
\section{Introduction and aim}

\paragraph{Aim.} We show that the canonical and covariant structures of Loop Quantum Gravity (LQG)---Ashtekar--Barbero variables, holonomy--flux representation and spin networks, quantum constraints, and spin--foam amplitudes---are obtained from the \emph{same} coherent fixed-point geometry that underlies all containments in MTT (modules G,S,$\Pi$,F,E).\footnote{See the Superset overview for modules and standing assumptions.} The Barbero--Immirzi parameter $\gBI$ is not free: it is mapped to a \emph{modal overlap ratio} fixed at the coherent point. This ties LQG predictions to the MTT bottleneck vector $\Theta$ and to the UV-finite QG form factor. % (Superset doc)

\paragraph{Context.} LQG quantizes real SU(2) connections and densitized triads on a spatial slice, producing a background-independent kinematics with spin-network basis and discrete geometric spectra, and covariant spin--foam dynamics \cite{Ashtekar1986,Barbero1995,Immirzi1997,Rovelli2004,Thiemann2007,AL1995,ALM1995,LOST2006,Perez2013,Engle2008,FreidelKrasnov2007}. We embed these ingredients in the MTT fixed-point framework developing the dictionary and constraints.

\paragraph{Standing modules.} We assume bounded geometry, a uniform spectral gap, a bounded joint projector $\Pi$ onto harmonic sectors, smoothing flow, and a convex energy functional near the fixed point (\S2).\footnote{These are exactly the modules used in your Superset paper.} 

\paragraph{Hypotheses and conventions.}
We work on a $3{+}1$ split $M_4=\mathbb{R}\times\Sigma$ of the coherent 4D face with:
(H1) $\Sigma$ is oriented, connected, piecewise-analytic, of bounded geometry (uniform curvature and injectivity bounds).
(H2) Non-degenerate tetrad and time gauge: choose an internal timelike unit $n^I$ so that $e^0{}_a=0$ and the internal gauge is reduced to $\mathrm{SU}(2)$.
(H3) The MTT 4D effective action contains Einstein--Hilbert and Holst-type terms with constant coefficients $\alpha_{\rm EH}(\Theta),\alpha_{\rm Holst}(\Theta)$ defined by modal overlap integrals on the coherent sector.
(H4) The projector $\Pi$ commutes with the natural action of analytic diffeomorphisms on the coherent fields, so the induced GNS state is diffeomorphism invariant.
All statements below are made under (G,S,$\Pi$,F,E) and (H1)--(H4).


% =========================
\section{Canonical variables from the MTT fixed point}

Let $M_4= \RR \times \Sigma$ be a $3+1$ foliation of the 4D face of the MTT coherent fixed point, with spatial slice $\Sigma$ of bounded geometry. Let $e^I_{\ a}$ be the spatial triad, $\omega^{IJ}_{\ a}$ the spin connection, and $K^i_{\ a}$ the extrinsic-curvature one-form in an internal $\mathfrak{su}(2)$ basis ($i=1,2,3$). 

\begin{defn}[MTT$\to$Ashtekar--Barbero dictionary]\label{def:AB}
In time gauge (H2), let $e^i{}_{a}$ be a spatial triad on $\Sigma$, $\Gamma^i{}_{a}(e)$ its torsion-free spin connection, and $K^i{}_{a}$ the extrinsic curvature one-form. Define
\[
A^i{}_{a}=\Gamma^i{}_{a}+\gamma_{\rm BI}\,K^i{}_{a},\qquad
E^{a}{}_{i}=\tfrac12\,\epsilon^{abc}\epsilon_{ijk}\,e^j{}_{b}\,e^k{}_{c}.
\]

Here $\gBI$ denotes the Barbero--Immirzi parameter (fixed later in \S\ref{subsec:ImmirziMap}).

The symplectic structure induced by the MTT action on the coherent sector yields the real bracket
\begin{equation}\label{eq:PB-AB}
\{A^{i}{}_{a}(x),\,E^{b}{}_{j}(y)\}
= 8\pi G\,\gamma_{\mathrm{BI}}\,\delta^{i}{}_{j}\,\delta^{b}{}_{a}\,\delta^{(3)}(x,y).
\end{equation}
which matches canonical LQG.
\end{defn}



\paragraph{Hypotheses for the canonical sector.}
We work on $(\Sigma,\bar q_{ab})$ of bounded geometry in the (piecewise) analytic category used in LOST, with
triads $e^i{}_a$ and torsion-free $\Gamma^i{}_a(e)$. The MTT 4D effective action on the coherent face contains
Einstein--Hilbert and Holst terms with $\Theta$–dependent coefficients,
\[
S_{\mathrm{eff}}[e,\omega] =
\alpha_{\mathrm{EH}}(\Theta)\int
\epsilon_{IJKL}\,e^{I}\wedge e^{J}\wedge F^{KL}[\omega]
\;+\;
\alpha_{\mathrm{Holst}}(\Theta)\int e^{I}\wedge e^{J}\wedge F_{IJ}[\omega].
\footnote{In standard conventions one writes
$S_{\mathrm{Holst}} \propto \tfrac{1}{\gamma}\int e\wedge e\wedge F[\omega]$.
Thus identifying $\alpha_{\mathrm{Holst}}(\Theta)\propto 1/\gamma$ yields
$\gamma_{\mathrm{MTT}}(\Theta)=\alpha_{\mathrm{EH}}(\Theta)/\alpha_{\mathrm{Holst}}(\Theta)$.}
\]
so that $\gamma_{\rm MTT}(\Theta):=\alpha_{\rm EH}(\Theta)/\alpha_{\rm Holst}(\Theta)$. Performing the $3{+}1$ split
and keeping the standard boundary term, the symplectic potential reduces in the coherent sector to
\[
\Theta_{\rm symp}
= \frac{1}{8\pi G}\int_\Sigma E^a{}_i\,\delta A^i{}_a,
\qquad
\{A^i{}_a(x),E^b{}_j(y)\}
= 8\pi G\,\gamma_{\rm MTT}(\Theta)\,\delta^i{}_j\,\delta^b{}_a\,\delta^{(3)}(x,y),
\]
i.e. with $\gamma=\gamma_{\rm MTT}(\Theta)$.

\begin{remark}[Holst and Nieh--Yan]
Under (H3) and vanishing torsion, the Holst density differs from a topological Nieh--Yan term; the ratio of the Einstein--Hilbert and Holst coefficients fixes the Immirzi parameter via \eqref{eq:gamma-map} below. This also reproduces \eqref{eq:PB-AB}.
\end{remark}


\begin{remark}[Holst term from MTT]
The MTT effective action on the 4D face contains, besides the Einstein--Hilbert term, a parity-odd piece of Holst type $S_{\mathrm{Holst}}\!\propto\!\gBI^{-1}\!\int e\wedge e \wedge F[\omega]$. The \emph{ratio} of these two terms is fixed by modal overlap integrals, giving a prediction $\gBI=\gBI_{\mathrm{MTT}}(\Theta)$ (see \S\ref{subsec:ImmirziMap}). \cite{Holst1996}
\end{remark}

% =========================
\section{Holonomy--flux algebra and the LOST representation}

Let $\Gamma$ be a piecewise analytic graph on $\Sigma$ with oriented edges $e$ and dual faces $S$. For $A\in\Acal$, define the holonomy $h_e[A]\in \mathrm{SU}(2)$ and flux $E_i(S)=\int_S \epsilon_{abc}\,E^a_{\ i}\,\dd x^b\wedge \dd x^c$.

\paragraph{LOST hypotheses.}
We work with piecewise analytic graphs on $\Sigma$; the holonomy--flux $*$-algebra generated by $\{h_e[A],E_i(S)\}$;
a diffeomorphism-invariant cyclic vector (the $\Pi$-induced vacuum); and continuity of parallel transport maps along
edges. These are precisely the assumptions under which the Ashtekar--Lewandowski (AL) representation is unique (LOST).

\begin{defn}[Holonomy--flux *-algebra]
Cylindrical functions on graphs $\gBI$ form $\Ccal_\gBI$; the holonomy--flux *-algebra is generated by $\{h_e, E_i(S)\}$ with relations induced by \eqref{eq:PB-AB}, Gauss rotations, and spatial diffeomorphisms.
\end{defn}

\begin{theorem}[MTT$\to$LOST]\label{thm:LOST}
Let $\mathfrak{A}$ be the holonomy--flux $*$-algebra built from holonomies along piecewise-analytic edges and fluxes through piecewise-analytic surfaces on $\Sigma$. Assume:
(i) $\mathrm{Diff}_{\mathrm{an}}(\Sigma)$ acts by automorphisms on $\mathfrak{A}$;
(ii) there exists a $\mathrm{Diff}^{\rm an}$-invariant cyclic state $\omega$; and
(iii) holonomies act continuously along analytic edges.
Then the GNS representation of $(\mathfrak{A},\omega)$ is unitarily equivalent to the Ashtekar--Lewandowski representation with the AL measure on generalized connections. In the MTT embedding, (H1) and (H4) furnish (i)--(iii), hence LOST uniqueness applies. \cite{AL1995,ALM1995,LOST2006}.
\end{theorem}


\begin{proof}[Idea]
$\Pi$ eliminates non-harmonic (KK-like) sectors and enforces background independence on $\Sigma$; the induced GNS construction satisfies the LOST axioms, hence uniqueness of the AL representation.
\end{proof}

\begin{cor}[Spin networks]
The AL kinematical Hilbert space is $L^2(\overline{\Acal},\dd\mu_{\mathrm{AL}})$ with orthonormal basis of spin networks $\psi_{\gBI,\{j_e\},\{\iota_v\}}$.
\end{cor}

% =========================
\section{Constraints and dynamics}

\paragraph{Constraints.} The Gauss, diffeomorphism, and Hamiltonian constraints descend from MTT gauge/diffeomorphism invariance and from the fixed-point Hamiltonian in the $3+1$ split:
\begin{align}
\mathcal{G}_i &= D_a E^a_{\ i} \approx 0,\qquad 
\mathcal{V}_a = E^b_{\ i} F^i_{\ ab} - (1+\gBI^2)\,K^i_{\ a}\,\mathcal{G}_i \approx 0,\\
\mathcal{H} &= \frac{E^a_{\ i} E^b_{\ j}}{\sqrt{\det E}}\left(\epsilon^{ij}\!_{k}\,F^k_{\ ab} - 2(1+\gBI^2)\,K^i_{\ [a}K^j_{\ b]}\right)\approx 0,
\end{align}
where $F^i_{\ ab}$ is the curvature of $A^i_{\ a}$.

\paragraph{Quantization.} In the AL representation the Gauss/diffeomorphism constraints are implemented by group-averaging; for the Hamiltonian we may use Thiemann's regularization or the Master-constraint \cite{Thiemann2007,ThiemannMaster2006}.

\paragraph{Domains and positivity.}
The Gauss and diffeomorphism constraints are implemented by group averaging on the cylindrical, diffeomorphism-invariant
dense domain. The Master-constraint $M=\int_\Sigma \frac{H^2}{\sqrt{\det E}}\,\mathrm{d}^3x$ defines a positive
quadratic form on the diffeo-invariant span of spin networks; closability and existence of a self-adjoint extension hold
under the standard assumptions (cylindrical consistency and graph-changing regularisation).

\begin{theorem}[Master-constraint in the MTT embedding]
There exists a positive, diffeomorphism-invariant Master-constraint $\mathbf{M}=\int_\Sigma \frac{\mathcal{H}^2}{\sqrt{\det E}}\dd^3x$ whose quadratic form is well-defined on diffeo-invariant states induced by the $\Pi$-projected sector. Its kernel coincides with the simultaneous solution space of $(\mathcal{G}_i,\mathcal{V}_a,\mathcal{H})$ in the embedding.
\end{theorem}

% =========================
\section{Geometric operators and the Immirzi map}
\label{subsec:ImmirziMap}

\begin{equation}
\label{eq:area-volume}
\hat A(S)\,\psi
= 8\pi\,\gamma_{\mathrm{BI}}\,\ell_{\mathrm P}^{2}\!
\sum_{e\cap S}\!\sqrt{j_{e}(j_{e}+1)}\,\psi,
\qquad
\hat V(R)\,\psi
= \sum_{v\in R}\sqrt{\hat Q_{v}}\,\psi,
\end{equation}
where $\widehat{Q}_v$ is the standard, graph-local, gauge-invariant node operator built from fluxes at $v$.
Both operators are essentially self-adjoint on the cylindrical domain.
\cite{RovelliSmolin1995,AshtekarLewandowski1997,Thiemann2007}.


\paragraph{Predictive $\gBI$.} In the MTT reduction, the 4D effective action contains the Einstein--Hilbert term and a Holst-like term with coefficients given by \emph{modal overlaps}:
\begin{equation}
S_{\mathrm{MTT}\to 4\mathrm{D}} \;\sim\; \alpha_{\mathrm{EH}}(\Theta)\,\int \epsilon_{IJKL}\,e^I\wedge e^J\wedge F^{KL}
\;+\; \alpha_{\mathrm{Holst}}(\Theta)\,\int e_I\wedge e_J\wedge F^{IJ} \,.
\end{equation}
Matching to the canonical symplectic structure gives
\begin{equation}
\label{eq:gammaMap}
\gBI \;=\; \gBI_{\mathrm{MTT}}(\Theta)\;:=\;\frac{\alpha_{\mathrm{EH}}(\Theta)}{\alpha_{\mathrm{Holst}}(\Theta)}\,.
\end{equation}

Matching the canonical symplectic form and \eqref{eq:PB-AB} fixes
\begin{equation}\label{eq:gamma-map}
\gBI_{\rm BI}=\gBI_{\rm BI}^{\rm MTT}(\Theta):=\frac{\alpha_{\rm EH}(\Theta)}{\alpha_{\rm Holst}(\Theta)}\,,
\end{equation}

so $\gamma_{\rm BI}$ is a \emph{derived} constant determined by the bottleneck vector $\Theta$ (spectral gaps, harmonic norms, volumes, curvature/overlap integrals). This directly links LQG geometric spectra and black-hole microstate counts to the same modal data that govern gauge/Yukawa/EFT parameters.

\begin{theorem}[Predictive Immirzi map]\label{thm:gammaMTT}
Under the standing MTT hypotheses (bounded geometry, $\Pi$-projection, coherent fixed point), the canonical symplectic
form induced by $S_{\rm eff}$ equals the Ashtekar--Barbero form with
\[
\gamma=\gamma_{\rm MTT}(\Theta)=\frac{\alpha_{\rm EH}(\Theta)}{\alpha_{\rm Holst}(\Theta)}.
\]
Consequently, the spectra of $\widehat{A}(S)$ and $\widehat{V}(R)$ are the standard LQG spectra with this value of
$\gamma$, and black-hole microstate counting constraints on $\gamma$ become constraints on~$\Theta$.
\end{theorem}


\begin{remark}[Cross-checks]
(i) Black-hole entropy fixes $\gBI$ in LQG microstate counts, giving a test of \eqref{eq:gammaMap}. 
(ii) Semiclassical weave/coherent states provide an independent constraint via area/volume calibration.
\end{remark}

% =========================
\section{Covariant dynamics from MTT: BF + simplicity $\Rightarrow$ EPRL/FK}

\paragraph{Holst/Plebanski form.} The MTT action in the 4D face admits a BF-like rewrite with simplicity constraints. Imposing simplicity with the $\gBI$-dependent linear constraints yields the EPRL/FK vertex amplitude on 2-complexes dual to triangulations \cite{Engle2008,FreidelKrasnov2007,Perez2013}.

We use the Holst/Plebanski rewrite with linear simplicity constraints adapted to $\gamma_{\rm BI}$;
imposing them weakly in the boundary Hilbert space yields EPRL/FK vertex amplitudes. In the large-spin asymptotics ($j\to\infty$ with areas fixed), stationary phase analysis reproduces Regge gravity with the correct Immirzi dependence. \cite{Barrett2009}

\paragraph{Simplicity constraints.}
We impose the linear simplicity constraints corresponding to $\gamma_{\rm MTT}(\Theta)$ so that $SL(2,\mathbb{C})$
representations $(\rho,k)$ map to $SU(2)$ spins $j$ via the standard $\gamma$-dependent embedding (e.g. $j=k$, $\rho=\gamma k$).
This fixes the EPRL/FK vertex labelling and ensures the correct Holst dependence in the large-spin asymptotics.

We employ the standard embedding for (unitary) principal series labels $(\rho,k)$ into $SU(2)$
spins via
\(
j=k,\quad \rho=\gamma_{\mathrm{BI}}\,k,
\)
so that boundary $SU(2)$ spins and intertwiners are $\gamma_{\mathrm{BI}}$–compatible with the Holst term.


\begin{theorem}[Spin--foam emergence]
The MTT path integral on the coherent sector reduces, after BF + simplicity reduction with the ratio \eqref{eq:gammaMap}, to a spin--foam sum with EPRL/FK vertex weights. In the large-spin limit the amplitude reproduces Regge gravity with the correct Barbero--Immirzi dependence.
\end{theorem}

\begin{remark}[Compatibility with the UV-finite QG sector]
The Stieltjes/Bernstein two-point structure and SPT Gaussian damping derived in the perturbative MTT QG sector persist in the background-independent sum as vertex/face amplitude suppressions at large momenta, consistent with cylindrical consistency and diffeomorphism invariance.
\end{remark}

% =========================
\section{Semiclassical sector and cosmology}

\paragraph{Complexifier coherent states.} Using Thiemann's complexifier construction adapted to the MTT symplectic form, one obtains semiclassical states peaked on $(A,E)$ reproducing classical geometry at scales $\gg \ell_P$ \cite{Thiemann2007}.

\paragraph{Weave states.} MTT-induced spin networks with typical edge spacing $\ell\sim \ell_P$ and label distribution set by $\Theta$ act as weaves approximating a given classical metric on $\Sigma$.

\paragraph{Loop quantum cosmology (outline).} The MTT $\to$ LQG map specializes to homogeneous/isotropic symmetry reduction, yielding effective Friedmann equations with a bounce at critical density $\rho_c\propto 1/\gBI^3$ times MTT scalings; thus $\rho_c$ is ultimately a function of $\Theta$ via \eqref{eq:gammaMap} (testable in early-universe scenarios). \cite{Ashtekar2006}


In standard LQC reductions one finds a critical density scaling $\rho_c\propto \gamma_{\rm BI}^{-3}$ (in natural units), so \eqref{eq:gamma-map} makes $\rho_c$ a function of $\Theta$. Coherent-state calibration of $\widehat{A}$ and $\widehat{V}$ then provides an independent semiclassical test of $\gamma_{\rm BI}^{\rm MTT}(\Theta)$.


% =========================
\section{Consistency with MTT’s perturbative QG finiteness}

\paragraph{Kinematics.} The AL representation and holonomy--flux algebra are consistent with OS positivity and causal support of TT two-point functions when probing around the coherent point (no conflict with background independence).

\paragraph{Dynamics.} SPT factorization implies an \emph{effective} Gaussian damping on internal graviton lines of Feynman-like expansions of spin--foam transition amplitudes, preserving cylindrical consistency while ensuring UV suppression in mixed representations.

\paragraph{Relation to overlaps.}
The same $\Theta$ that fixes $\gamma_{\rm BI}$ in \eqref{eq:gamma-map} appears throughout the Superset overlaps (Planck mass, gauge couplings, Yukawas, EFT/KK scales), so LQG geometric spectra are predicted jointly with non-gravitational observables; see the Superset discussion of bottlenecks and global fits. 

\begin{lemma}[Gaussian suppression in mixed representations]
In the $\Pi$-projected coherent sector, internal graviton propagators entering Feynman-like expansions of spin-foam
transition amplitudes carry an entire/Gaussian form factor with scale set by the MTT spectral gap. Consequently, large-momentum
(or large-spin with fixed geometry) sectors are exponentially suppressed, compatibly with cylindrical consistency.
\end{lemma}

\paragraph{Coherence with the perturbative sector.}
The Stieltjes/Bernstein positivity of TT two-point functions, causal support properties, and the all-orders BV/QME
renormalisation established in the perturbative MTT QG construction carry over here at the level of kinematics and regulator removal;
see the companion QG paper for full details.

% =========================
\section{Conclusions and tests}

We have constructed a direct MTT$\to$LQG embedding, fixed the Barbero--Immirzi parameter by modal overlaps, and matched both canonical and covariant LQG structures. Key tests:
\begin{itemize}
\item \textbf{Black-hole entropy:} Check that $\gBI_{\mathrm{MTT}}(\Theta)$ matches the value required by microstate counting within uncertainties.
\item \textbf{Semiclassical calibration:} Compare area/volume expectation values on weave/coherent states to classical geometry induced by the MTT fixed point.
\item \textbf{Cosmology:} Use $\gBI_{\mathrm{MTT}}(\Theta)$ in LQC phenomenology (e.g.\ $\rho_c$) and confront with data.
\item \textbf{Amplitudes:} Verify Regge asymptotics and form-factor suppression scale derived from the SPT map.
\end{itemize}

\paragraph{Outlook.} Extend the EPRL/FK derivation to include matter (gauge and fermions) from the same modal overlaps; quantify the relation between $\gBI_{\mathrm{MTT}}$ and the UV-finite QG form factor at low spin; develop global fits in which $\gBI$ is predicted jointly with gauge/Yukawa parameters.

\bigskip

\bibliographystyle{plain} % or abbrv, unsrt, alpha, etc
\bibliography{main}

\end{document}
\typeout{get arXiv to do 4 passes: Label(s) may have changed. Rerun}
